% Generated draft: requires mathematical and source review.
\documentclass[11pt]{article}
\usepackage[margin=25mm]{geometry}
\usepackage{fontspec}
\setmainfont{Latin Modern Roman}
\usepackage{amsmath,amssymb,mathtools,unicode-math}
\setmathfont{Latin Modern Math}
\usepackage{enumitem,needspace,xurl,hyperref}
\setlength{\parindent}{0pt}
\setlength{\parskip}{5pt}
\setlength{\emergencystretch}{3em}
\newcommand{\N}{\mathbb{N}}
\newcommand{\Z}{\mathbb{Z}}
\newcommand{\Q}{\mathbb{Q}}
\newcommand{\R}{\mathbb{R}}
\newcommand{\C}{\mathbb{C}}
\newcommand{\F}{\mathbb{F}}
\newcommand{\A}{\mathbb{A}}
\newcommand{\PP}{\mathbb{P}}
\newcommand{\E}{\mathbb{E}}
\newcommand{\eps}{\varepsilon}
\newcommand{\dd}{\,\mathrm{d}}
\begin{document}
\section*{UnsolvedMath ID: 30002871}
\textbf{Source problem code:} catalogue entry
\subsection{Definitions and mathematical statement}
Write the intended integer column configuration as \(A=(a_1,\ldots,a_m)\), with \(a_j\in\Z^{d+1}\), solely to explain the semigroup notation in the record. Here \(m\) denotes the number of columns; the record does not identify it with its symbol \(n\). Define
\[
\N A=\left\{\sum_{j=1}^{m}u_j a_j : u_j\in\Z_{\geq 0}\right\}\subseteq\Z^{d+1}.
\]
For \(c=(c_0,\ldots,c_d)\in\N A\), the integer \(c_0\) is its first component. The notation \(M_A^\beta\) denotes the GKZ module with configuration \(A\) and parameter \(\beta\), as named in the source. Its ambient space, defining convention, and the admissible configurations \(A\) are not specified in the supplied record.

An increasing filtration on a module \(M\) is a family \((F_pM)_{p\in\Z}\) with \(F_pM\subseteq F_{p+1}M\). The source uses \(F^H\) for a Hodge filtration and \(F^{\mathrm{ord}}\) for an order filtration. In a specified cyclic differential-operator presentation \(M=D/I\), the usual order filtration is the image of operators of order at most \(p\) in \(M\), with zero terms for \(p<0\). The record supplies neither such a presentation and its normalization nor a Hodge-module realization determining \(F^H\).

\textbf{Recorded assertion, subject to unresolved hypotheses.} For the intended configurations and parameters \(c\in\N A\), the source asserts, for every \(p\in\Z\),
\[
F^H_p M_A^0=F^{\mathrm{ord}}_{p+d}M_A^0,
\qquad
F^H_p M_A^{-c}=F^{\mathrm{ord}}_{p-n-c_0}M_A^{-c}.
\]
These are equalities of filtration terms on the respective modules. The symbol \(n\), the class of admissible configurations, and any additional restrictions on \(c\) remain unspecified. In particular, rank, homogeneity, normality or saturation hypotheses cannot be inferred from the displayed assertion. Consequently, this is a faithful transcription of the proposed identities, not a fully specified universally quantified conjecture.

The supplied background instead discusses a GKZ \emph{weight}-filtration conjecture and a full saturated-semigroup assertion. It gives no formulation of that assertion or explanation connecting it to these \emph{Hodge}-filtration identities. That background qualification is retained as unresolved.
\subsection{Short English statement}
The recorded claim says that the Hodge filtrations on \(M_A^0\) and \(M_A^{-c}\) agree with their order filtrations after shifts by \(d\) and \(-n-c_0\), respectively, where \(c\in\N A\). The allowed configurations, the meaning of \(n\), and the Hodge structures are unspecified. The supplied partial-solution assessment concerns weight filtrations, whose connection to these identities requires review.
\subsection{Sources}
\begin{itemize}
\item Supplied local UnsolvedMath registry record, permanent ID \(30002871\), source problem code \texttt{catalogue entry}, title \emph{Hodge and Order Filtrations on GKZ Modules}; project registry \texttt{data/unsolvedmath/problems.json}. This attribution uses the supplied record; the local registry file was not independently inspected.
\item Supplied source citation: \emph{Mirror Symmetry, Hodge Theory and Differential Equations} (2016), Oberwolfach Reports, DOI \texttt{10.4171/owr/2015/22}, pp.~1201--1254. \url{https://doi.org/10.4171/owr/2015/22}. The literature-source list labels this reference ``Oberwolfach Report 22.'' The record assigns proposed year 2015 and supplies no named proposer.
\item Supplied literature reference: Reichelt--Walther, \emph{Weight Filtrations on GKZ-Systems}. \url{https://arxiv.org/abs/1809.04247}.
\item Literature assessment dated exactly \texttt{2026-08-21}: supplied dated curation, not independently verified. It labels the problem partially solved and states that the GKZ weight-filtration conjecture is established in normal and several low-dimensional or simplicial settings, with explicit image and filtration results, while the full saturated-semigroup assertion in the record is not proved generally. The supplied background describes this as literature triage derived from OWR curation and requests independent verification of time-sensitive status and cited claims. Its relationship to the displayed Hodge-filtration identities is unresolved.
\item The registry labels the statement \texttt{corrected\_verified} and reports manual semantic, typographic, and LaTeX review of a meaning-preserving cleanup. These are supplied editorial labels and do not establish the mathematical identities or resolve their missing hypotheses.
\end{itemize}
\end{document}
